\documentclass[12pt,a4paper]{article}

% XeLaTeX用のシンプルな日本語設定
\usepackage{fontspec}
\setmainfont{Noto Serif CJK JP}[Scale=1.0]
\setsansfont{Noto Sans CJK JP}[Scale=1.0]
\setmonofont{Noto Sans Mono CJK JP}[Scale=0.9]

% 数学パッケージ
\usepackage{amsmath,amssymb,amsthm}
\usepackage{mathtools}
\usepackage{tikz}
\usepackage{tikz-cd}
\usetikzlibrary{arrows,positioning,shapes,calc}

% その他のパッケージ
\usepackage{hyperref}
\usepackage{enumerate}
\usepackage{framed}
\usepackage{listings}
\usepackage{xcolor}

% ページ設定
\usepackage[top=25mm,bottom=25mm,left=25mm,right=25mm]{geometry}

% Lean コードの設定
\lstdefinelanguage{Lean}{
  keywords={import, def, theorem, lemma, example, structure, inductive, where, by, intro, exact, have, show, sorry, noncomputable, abbrev, section, end, namespace, variable, open},
  keywordstyle=\color{blue}\bfseries,
  sensitive=true,
  comment=[l]{--},
  commentstyle=\color{gray}\itshape,
  string=[b]",
  stringstyle=\color{red},
}

\lstset{
  language=Lean,
  basicstyle=\small\ttfamily,
  breaklines=true,
  frame=single,
  numbers=left,
  numberstyle=\tiny\color{gray},
}

% 定理環境
\theoremstyle{definition}
\newtheorem{definition}{定義}[section]
\newtheorem{theorem}[definition]{定理}
\newtheorem{lemma}[definition]{補題}
\newtheorem{proposition}[definition]{命題}
\newtheorem{example}[definition]{例}
\newtheorem{remark}[definition]{注意}
\newtheorem{corollary}[definition]{系}

% カスタムコマンド
\newcommand{\N}{\mathbb{N}}
\newcommand{\R}{\mathbb{R}}
\newcommand{\cF}{\mathcal{F}}
\newcommand{\sA}{\mathfrak{A}}
\DeclareMathOperator{\generateFrom}{generateFrom}
\DeclareMathOperator{\minLayer}{minLayer}

\title{\Large 確率論における構造塔\\
SigmaAlgebraTower \& StoppingTime: \\
フィルトレーションから停止時間へ}

\author{学部生向け数学解説\\
\vspace{5mm}
\normalsize Claude Code により生成\\
\normalsize Lean 4 コード: CODEX (OpenAI)\\
\normalsize \TeX 解説: Claude Code\\
\vspace{3mm}
\normalsize プロジェクト: Structure Tower Theory}

\date{2025年11月19日}

\begin{document}

\maketitle

\begin{abstract}
本稿では、構造塔理論を確率論に応用した形式化について解説する。特に、$\sigma$-代数の階層構造としてのフィルトレーション、停止時間と $\minLayer$ 関数の関係、停止過程の構成を扱う。\texttt{SigmaAlgebraTower\_Skeleton.lean} と \texttt{StoppingTime\_MinLayer.lean} の Lean 4 コードを基に、確率論における構造塔の理論的基礎と具体的実装を学部3-4年生向けに解説する。測度論と確率論の基礎知識を前提とする。
\end{abstract}

\tableofcontents

\section{序論：確率論における階層構造}

\subsection{フィルトレーションとは}

確率論において、\textbf{フィルトレーション}（filtration）は情報の増加を数学的に表現する概念である。

\begin{definition}[フィルトレーション]
可測空間 $(\Omega, \cF)$ 上の\textbf{フィルトレーション} $(\cF_n)_{n \in \N}$ とは、$\sigma$-代数の増加列
\[
\cF_0 \subseteq \cF_1 \subseteq \cF_2 \subseteq \cdots \subseteq \cF
\]
である。直感的には、$\cF_n$ は「時刻 $n$ までに得られた情報」を表す。
\end{definition}

\begin{center}
\begin{tikzpicture}[>=stealth]
  \node[draw,rectangle,minimum width=1.5cm,minimum height=1cm] (F0) at (0,0) {$\cF_0$};
  \node[draw,rectangle,minimum width=2cm,minimum height=1.5cm] (F1) at (2.5,0) {$\cF_1$};
  \node[draw,rectangle,minimum width=2.5cm,minimum height=2cm] (F2) at (5.5,0) {$\cF_2$};
  \node[draw,rectangle,minimum width=3cm,minimum height=2.5cm] (F3) at (9,0) {$\cF_3$};
  
  \draw[->,thick] (F0) -- (F1) node[midway,above] {包含};
  \draw[->,thick] (F1) -- (F2);
  \draw[->,thick] (F2) -- (F3);
  
  \node[above] at (0,0.8) {\small 時刻0の情報};
  \node[above] at (2.5,1.1) {\small 時刻1の情報};
  \node[above] at (5.5,1.4) {\small 時刻2の情報};
  \node[above] at (9,1.7) {\small 時刻3の情報};
\end{tikzpicture}
\end{center}

\subsection{構造塔との対応}

フィルトレーションは、まさに構造塔の典型例である：

\begin{itemize}
\item \textbf{基礎集合}：$\Omega$ の部分集合全体
\item \textbf{添字集合}：時刻 $\N$（自然な順序付き）
\item \textbf{層}：各時刻 $n$ での可測集合族 $\cF_n$
\item \textbf{単調性}：$n \leq m \Rightarrow \cF_n \subseteq \cF_m$
\end{itemize}

\subsection{本稿の構成}

\begin{enumerate}
\item 第\ref{sec:sigma}節：$\sigma$-代数の構造塔
\item 第\ref{sec:stopping}節：停止時間と $\minLayer$
\item 第\ref{sec:stopped}節：停止過程の構成
\item 第\ref{sec:applications}節：応用と展望
\end{enumerate}

\section{$\sigma$-代数の構造塔}\label{sec:sigma}

\subsection{$\sigma$-代数の順序構造}

\begin{definition}[$\sigma$-代数の順序]
可測空間 $\Omega$ 上の二つの $\sigma$-代数 $\cF_1, \cF_2$ に対して、
\[
\cF_1 \leq \cF_2 \overset{\text{def}}{\Longleftrightarrow} \cF_1 \subseteq \cF_2
\]
と定義する。すなわち、$\cF_1$ で可測な集合は全て $\cF_2$ でも可測。
\end{definition}

この順序により、$\sigma$-代数全体は半順序集合をなす。

\begin{remark}[粗いと細かい]
$\cF_1 \leq \cF_2$ のとき：
\begin{itemize}
\item $\cF_1$ は\textbf{粗い}（coarser）：区別できる事象が少ない
\item $\cF_2$ は\textbf{細かい}（finer）：区別できる事象が多い
\end{itemize}
\end{remark}

\subsection{SigmaAlgebraTower の定義}

\begin{definition}[σ-代数の構造塔]
$\Omega$ を可測空間とする。$\sigma$-代数を添字とする構造塔を次のように定義する：
\begin{itemize}
\item \textbf{基礎集合}：$\mathcal{P}(\Omega)$（$\Omega$ の部分集合全体）
\item \textbf{添字集合}：$\text{MeasurableSpace}(\Omega)$（$\Omega$ 上の $\sigma$-代数全体）
\item \textbf{層 $\cF$}：$\{A \subseteq \Omega \mid A \text{ は } \cF \text{ で可測}\}$
\item \textbf{minLayer}：各集合 $A$ に対して、$A$ を可測にする最小の $\sigma$-代数 $\generateFrom\{A\}$
\end{itemize}
\end{definition}

Lean 4 での実装（\texttt{SigmaAlgebraTower\_Skeleton.lean} 71-108行目）：

\begin{lstlisting}[language=Lean]
def SigmaAlgebraTower : StructureTowerMin (Set Ω) (MeasurableSpace Ω) where
  layer 𝓕 := {A : Set Ω | @MeasurableSet Ω 𝓕 A}
  
  covering := by
    intro A
    use ⊤  -- 全体のσ-代数
    trivial
  
  monotone := by
    intro 𝓕₁ 𝓕₂ h𝓕 A hA
    exact h𝓕 A hA
  
  minLayer := fun A =>
    MeasurableSpace.generateFrom {A}
  
  minLayer_mem := by
    intro A
    exact MeasurableSpace.measurableSet_generateFrom
      (by simp : A ∈ ({A} : Set (Set Ω)))
  
  minLayer_minimal := by
    intro A 𝓕 hA
    apply MeasurableSpace.generateFrom_le
    intro B hB
    rcases hB with rfl
    exact hA
\end{lstlisting}

\subsection{生成 $\sigma$-代数の最小性}

\begin{theorem}[generateFrom の最小性]
集合族 $S \subseteq \mathcal{P}(\Omega)$ と $\sigma$-代数 $\cF$ に対して、
\[
(\forall B \in S,\ B \in \cF) \Longrightarrow \generateFrom(S) \subseteq \cF
\]
が成り立つ。すなわち、$\generateFrom(S)$ は $S$ を含む最小の $\sigma$-代数である。
\end{theorem}

\begin{center}
\begin{tikzpicture}[>=stealth]
  \node[draw,circle,minimum size=3cm] (gen) at (0,0) {};
  \node[draw,circle,minimum size=5cm] (F) at (0,0) {};
  
  \node at (0,-0.5) {$\generateFrom(S)$};
  \node at (0,2) {$\cF$};
  
  \fill[red] (-0.3,0.3) circle (2pt) node[left] {$B_1 \in S$};
  \fill[red] (0.4,-0.2) circle (2pt) node[right] {$B_2 \in S$};
  
  \draw[->,thick,blue] (-1.2,-1.2) -- (-2,-2) node[below] {最小の $\sigma$-代数};
\end{tikzpicture}
\end{center}

\subsection{構造塔としての性質}

\begin{proposition}[単調性]
$\cF_1 \leq \cF_2$ ならば、層の包含関係
\[
\text{layer}(\cF_1) \subseteq \text{layer}(\cF_2)
\]
が成り立つ。すなわち、$\cF_1$-可測な集合は $\cF_2$-可測である。
\end{proposition}

\begin{proposition}[被覆性]
任意の集合 $A \subseteq \Omega$ に対して、ある $\sigma$-代数 $\cF$（例えば $\cF = \top$）が存在して、
\[
A \in \text{layer}(\cF)
\]
が成り立つ。
\end{proposition}

\section{フィルトレーションと構造塔}\label{sec:filtration}

\subsection{フィルトレーションの形式化}

\begin{definition}[離散フィルトレーション]
可測空間 $(\Omega, \cF)$ 上の\textbf{離散フィルトレーション}とは、単調な $\sigma$-代数の族
\[
(\cF_n)_{n \in \N} : \N \to \text{MeasurableSpace}(\Omega)
\]
であって、$n \leq m \Rightarrow \cF_n \subseteq \cF_m$ を満たすものである。
\end{definition}

Lean 4 での実装：

\begin{lstlisting}[language=Lean]
structure Core where
  𝓕 : ι → MeasurableSpace Ω
  mono : Monotone 𝓕
\end{lstlisting}

\subsection{被覆性を持つフィルトレーション}

確率論の応用のため、「すべての集合がいずれ可測になる」という性質を追加する。

\begin{definition}[被覆性を持つフィルトレーション]
フィルトレーション $(\cF_n)_{n \in \N}$ が\textbf{被覆性}（covering property）を持つとは、
\[
\forall A \subseteq \Omega,\ \exists n \in \N,\ A \in \cF_n
\]
が成り立つことをいう。
\end{definition}

\begin{lstlisting}[language=Lean]
structure SigmaAlgebraFiltrationWithCovers where
  base : SigmaAlgebraFiltration.Core (Ω := Ω) (ι := ℕ)
  covers : ∀ A : Set Ω, ∃ n : ℕ, @MeasurableSet Ω (base.𝓕 n) A
\end{lstlisting}

\subsection{フィルトレーションから構造塔へ}

被覆性を持つフィルトレーション $\cF$ から、構造塔 $\text{TowerOf}(\cF)$ を構成できる：

\begin{theorem}[フィルトレーションの構造塔化]
被覆性を持つフィルトレーション $\cF = (\cF_n)_{n \in \N}$ に対して、構造塔
\[
\text{TowerOf}(\cF) : \text{StructureTowerMin}(\mathcal{P}(\Omega), \N)
\]
が次のように定義される：
\begin{itemize}
\item $\text{layer}(n) = \{A \subseteq \Omega \mid A \in \cF_n\}$
\item $\minLayer(A) = \min\{n \in \N \mid A \in \cF_n\}$（Nat.find を使用）
\end{itemize}
\end{theorem}

\begin{proof}
被覆性より、各 $A$ に対して $A \in \cF_n$ となる $n$ が存在する。Nat.find により最小の $n$ が一意に定まる。単調性と最小性は定義から従う。
\end{proof}

\section{停止時間と $\minLayer$}\label{sec:stopping}

\subsection{停止時間の定義}

\begin{definition}[停止時間]
フィルトレーション $(\cF_n)_{n \in \N}$ に対して、確率変数 $\tau : \Omega \to \N$ が\textbf{停止時間}（stopping time）であるとは、
\[
\forall n \in \N,\ \{\omega \in \Omega \mid \tau(\omega) \leq n\} \in \cF_n
\]
が成り立つことをいう。
\end{definition}

\begin{remark}[停止時間の直感]
停止時間 $\tau(\omega)$ は、「標本点 $\omega$ に関する情報が初めて明らかになる時刻」を表す。重要なのは、$\tau(\omega) \leq n$ かどうかを判定するのに、時刻 $n$ までの情報 $\cF_n$ で十分であることである。
\end{remark}

\begin{center}
\begin{tikzpicture}[>=stealth,scale=0.8]
  % 時間軸
  \draw[->,thick] (0,0) -- (10,0) node[right] {時刻 $n$};
  \foreach \x in {0,2,4,6,8}
    \draw (\x,0.1) -- (\x,-0.1) node[below] {$\x$};
  
  % 標本点
  \draw[thick,blue] (0,3) node[left] {$\omega_1$} -- (4,3);
  \fill[blue] (4,3) circle (3pt) node[above] {$\tau(\omega_1)=4$};
  
  \draw[thick,red] (0,2) node[left] {$\omega_2$} -- (6,2);
  \fill[red] (6,2) circle (3pt) node[above] {$\tau(\omega_2)=6$};
  
  % 停止集合
  \fill[yellow,opacity=0.3] (0,0) rectangle (4,3.5);
  \node at (2,3.5) {$\{\tau \leq 4\} \in \cF_4$};
\end{tikzpicture}
\end{center}

\subsection{停止集合と構造塔}

\begin{theorem}[停止集合の層への所属]
停止時間 $\tau$ に対して、各 $n \in \N$ で
\[
\{\omega \mid \tau(\omega) \leq n\} \in \text{layer}(n)
\]
が成り立つ。すなわち、停止集合は対応する層に属する。
\end{theorem}

\begin{proof}
停止時間の定義より、$\{\tau \leq n\} \in \cF_n$ である。構造塔の層 $\text{layer}(n)$ は $\cF_n$ の可測集合全体なので、$\{\tau \leq n\} \in \text{layer}(n)$ が成り立つ。
\end{proof}

Lean 4 での実装：

\begin{lstlisting}[language=Lean]
lemma stopping_sets_in_tower (ℱ : Filtration Ω)
    (τ : StoppingTime ℱ) (n : ℕ) :
    {ω : Ω | τ.τ ω ≤ n} ∈ (towerOf (Ω := Ω) ℱ).layer n := by
  change @MeasurableSet Ω (ℱ.base.𝓕 n) {ω : Ω | τ.τ ω ≤ n}
  exact τ.measurable n
\end{lstlisting}

\subsection{単点の初可測時刻}

構造塔の $\minLayer$ を使って、新しい概念を導入できる。

\begin{definition}[初可測時刻]
標本点 $\omega \in \Omega$ に対して、\textbf{初可測時刻}（first measurable time）を
\[
\text{firstMeasurableTime}(\omega) := \minLayer_{\text{TowerOf}(\cF)}(\{\omega\})
\]
と定義する。これは、単点集合 $\{\omega\}$ が初めて可測になる時刻である。
\end{definition}

\begin{theorem}[初可測時刻での可測性]
任意の $\omega \in \Omega$ に対して、
\[
\{\omega\} \in \cF_{\text{firstMeasurableTime}(\omega)}
\]
が成り立つ。
\end{theorem}

\begin{theorem}[初可測時刻の最小性]
$\{\omega\} \in \cF_n$ ならば、$\text{firstMeasurableTime}(\omega) \leq n$ である。
\end{theorem}

\subsection{停止時間と $\minLayer$ の関係}

停止時間は、各標本点の「初可測時刻」と密接に関連している。

\begin{remark}[概念的対応]
停止時間 $\tau$ は、以下のように解釈できる：
\begin{itemize}
\item $\tau(\omega)$ は、$\omega$ に関する「何らかの情報」が初めて分かる時刻
\item これは、$\omega$ を含む事象が初めて可測になる時刻と対応
\item 構造塔の言葉では、$\minLayer$ の概念そのもの
\end{itemize}
\end{remark}

\begin{center}
\begin{tikzpicture}[>=stealth]
  \node[draw,rectangle,minimum width=3cm,minimum height=1.5cm] (stop) at (0,0) {停止時間 $\tau(\omega)$};
  \node[draw,rectangle,minimum width=3cm,minimum height=1.5cm] (min) at (6,0) {$\minLayer(\{\omega\})$};
  
  \draw[->,thick,blue] (stop) -- (min) node[midway,above] {構造塔の視点};
  
  \node[below=0.5cm of stop,text width=3.5cm,align=center] {\small $\omega$ に関する情報が\\初めて分かる時刻};
  \node[below=0.5cm of min,text width=3.5cm,align=center] {\small $\{\omega\}$ が初めて\\可測になる時刻};
\end{tikzpicture}
\end{center}

\section{停止過程の構成}\label{sec:stopped}

\subsection{停止過程の定義}

\begin{definition}[停止過程]
確率過程 $X = (X_n)_{n \in \N}$ と停止時間 $\tau$ に対して、\textbf{停止過程}（stopped process）$X^\tau$ を
\[
X^\tau_n(\omega) := X_{\min(n, \tau(\omega))}(\omega)
\]
と定義する。すなわち、停止時刻 $\tau(\omega)$ 以降は値が固定される。
\end{definition}

\begin{center}
\begin{tikzpicture}[>=stealth,scale=0.8]
  % 軸
  \draw[->,thick] (0,0) -- (10,0) node[right] {時刻 $n$};
  \draw[->,thick] (0,0) -- (0,4) node[above] {$X_n$};
  
  % 元の過程
  \draw[thick,blue] (0,1) -- (1,1.5) -- (2,1.8) -- (3,2.5) -- (4,2.2) -- (5,2.8) -- (6,3.2) -- (7,3.5) -- (8,3.3) -- (9,3.8);
  \node[blue] at (9,4.2) {元の過程 $X_n$};
  
  % 停止時刻
  \draw[red,dashed,thick] (4,0) -- (4,4);
  \node[red] at (4,-0.5) {$\tau(\omega)=4$};
  
  % 停止過程
  \draw[thick,green!60!black,line width=1.5pt] (0,1) -- (1,1.5) -- (2,1.8) -- (3,2.5) -- (4,2.2) -- (5,2.2) -- (6,2.2) -- (7,2.2) -- (8,2.2) -- (9,2.2);
  \node[green!60!black] at (7,1.5) {停止過程 $X^\tau_n$};
  
  \fill[red] (4,2.2) circle (4pt);
\end{tikzpicture}
\end{center}

Lean 4 での定義：

\begin{lstlisting}[language=Lean]
def stopped (X : ℕ → Ω → α) (τ : Ω → ℕ) (n : ℕ) : Ω → α :=
  fun ω => X (Nat.min n (τ ω)) ω
\end{lstlisting}

\subsection{停止 $\sigma$-代数}

\begin{definition}[停止 $\sigma$-代数]
停止時間 $\tau$ に対して、\textbf{停止 $\sigma$-代数} $\cF_\tau$ を次のように定義する：
\[
A \in \cF_\tau \overset{\text{def}}{\Longleftrightarrow} \forall n \in \N,\ A \cap \{\tau \leq n\} \in \cF_n
\]
\end{definition}

\begin{remark}[停止 $\sigma$-代数の直感]
$\cF_\tau$ は「停止時刻 $\tau$ までに得られる情報」を表す。事象 $A$ が $\cF_\tau$ に属するとは、各時刻 $n$ で「$A$ かつ $\tau \leq n$」が判定可能であることを意味する。
\end{remark}

\subsection{停止過程の可測性}

\begin{theorem}[停止過程の適合性]
フィルトレーション $(\cF_n)_{n \in \N}$ に適合する確率過程 $X$ と停止時間 $\tau$ に対して、停止過程 $X^\tau$ も $(\cF_n)_{n \in \N}$ に適合する。すなわち、
\[
\forall n \in \N,\ X^\tau_n \text{ は } \cF_n\text{-可測}
\]
が成り立つ。
\end{theorem}

\begin{proof}[証明の概要]
$X^\tau_n(\omega)$ の値は、$\{\tau \leq k\}$ と $\{\tau > k\}$ で場合分けされる。各場合で対応する $X_k$ が $\cF_n$-可測であることを示す。詳細は \texttt{StoppingTime\_MinLayer.lean} の 617-667行目を参照。
\end{proof}

\subsection{切り詰め停止時間}

\begin{definition}[切り詰め停止時間]
停止時間 $\tau$ と自然数 $n$ に対して、\textbf{切り詰め停止時間} $\tau \wedge n$ を
\[
(\tau \wedge n)(\omega) := \min(\tau(\omega), n)
\]
と定義する。
\end{definition}

\begin{proposition}[切り詰めも停止時間]
$\tau$ が停止時間ならば、$\tau \wedge n$ も停止時間である。
\end{proposition}

\begin{center}
\begin{tikzpicture}[>=stealth,scale=0.8]
  \draw[->,thick] (0,0) -- (10,0) node[right] {時刻};
  
  % 元の停止時間
  \fill[blue] (7,2) circle (3pt);
  \draw[blue,thick] (0,2) -- (7,2);
  \node[blue,above] at (7,2) {$\tau(\omega)=7$};
  
  % 切り詰め
  \draw[red,dashed,thick] (4,0) -- (4,3);
  \node[red] at (4,-0.5) {$n=4$};
  
  % 切り詰め後
  \fill[green!60!black] (4,1) circle (3pt);
  \draw[green!60!black,thick,line width=1.5pt] (0,1) -- (4,1);
  \node[green!60!black,below] at (4,1) {$\tau \wedge n = 4$};
  
  \draw[->,blue,thick] (7,1.8) to[bend left] (4.5,1.2);
\end{tikzpicture}
\end{center}

\section{応用と理論的意義}\label{sec:applications}

\subsection{オプショナル停止定理への準備}

停止過程の理論は、確率論の重要な結果である\textbf{オプショナル停止定理}（optional stopping theorem）の基礎となる。

\begin{theorem}[オプショナル停止定理（簡易版）]
マルチンゲール $(M_n, \cF_n)_{n \in \N}$ と有界停止時間 $\tau$ に対して、停止過程 $M^\tau$ もマルチンゲールである。
\end{theorem}

この定理の証明において、停止過程の可測性が本質的な役割を果たす。

\subsection{構造塔理論の視点}

構造塔の視点から見ると、確率論の多くの概念が統一的に理解できる：

\begin{center}
\begin{tikzpicture}[>=stealth]
  \node[draw,ellipse,minimum width=3cm,minimum height=1cm] (filt) at (0,3) {フィルトレーション};
  \node[draw,ellipse,minimum width=3cm,minimum height=1cm] (tower) at (0,0) {構造塔};
  \node[draw,ellipse,minimum width=3cm,minimum height=1cm] (stop) at (5,3) {停止時間};
  \node[draw,ellipse,minimum width=3cm,minimum height=1cm] (min) at (5,0) {$\minLayer$};
  
  \draw[->,thick,blue] (filt) -- (tower) node[midway,left] {TowerOf};
  \draw[->,thick,red] (stop) -- (min) node[midway,right] {対応};
  \draw[<->,thick,dashed] (filt) -- (stop);
  \draw[<->,thick,dashed] (tower) -- (min);
\end{tikzpicture}
\end{center}

\begin{itemize}
\item \textbf{フィルトレーション} $\longleftrightarrow$ \textbf{構造塔}：情報の増加 = 層の包含
\item \textbf{停止時間} $\longleftrightarrow$ \textbf{minLayer}：初可測時刻
\item \textbf{停止 $\sigma$-代数} $\longleftrightarrow$ \textbf{層の交叉}：停止時刻での情報
\item \textbf{適合過程} $\longleftrightarrow$ \textbf{層保存写像}：層の構造を保つ
\end{itemize}

\subsection{測度論的拡張}

現在の形式化は $\sigma$-代数レベルにとどまっているが、次のステップとして測度の導入が考えられる：

\begin{enumerate}
\item \textbf{測度の構造塔}：$(\Omega, \cF_n, \mu)$ の階層
\item \textbf{可積分性の伝播}：層間での可積分性の関係
\item \textbf{条件付き期待値}：$\mathbb{E}[X \mid \cF_n]$ と $\minLayer$ の関係
\item \textbf{マルチンゲール}：適合過程の特別な場合
\end{enumerate}

\subsection{他の数学分野への応用}

構造塔理論の枠組みは、確率論を超えて広がる可能性を持つ：

\begin{itemize}
\item \textbf{代数幾何学}：層の理論との対応
\item \textbf{トポロジー}：開集合の階層構造
\item \textbf{論理学}：知識の増加のモデル化
\item \textbf{計算機科学}：プログラムの状態遷移
\end{itemize}

\section{まとめ}

\subsection{本稿の成果}

本稿では、構造塔理論を確率論に応用した以下の内容を解説した：

\begin{enumerate}
\item $\sigma$-代数を構造塔として形式化
\item フィルトレーションの構造塔的解釈
\item 停止時間と $\minLayer$ の関係
\item 停止過程の構成と性質
\end{enumerate}

\subsection{形式化の階層}

\begin{center}
\begin{tikzpicture}[>=stealth]
  \node[draw,rectangle,minimum width=4cm,minimum height=1cm,fill=blue!10] (level1) at (0,0) 
    {Level 1: 構造塔（抽象）};
  \node[draw,rectangle,minimum width=4cm,minimum height=1cm,fill=blue!20] (level2) at (0,1.5) 
    {Level 2: $\sigma$-代数の構造塔};
  \node[draw,rectangle,minimum width=4cm,minimum height=1cm,fill=blue!30] (level3) at (0,3) 
    {Level 3: フィルトレーション};
  \node[draw,rectangle,minimum width=4cm,minimum height=1cm,fill=blue!40] (level4) at (0,4.5) 
    {Level 4: 停止時間と停止過程};
  \node[draw,rectangle,minimum width=4cm,minimum height=1cm,fill=blue!50] (level5) at (0,6) 
    {Level 5: マルチンゲール理論};
  
  \draw[->,thick] (level1) -- (level2);
  \draw[->,thick] (level2) -- (level3);
  \draw[->,thick] (level3) -- (level4);
  \draw[->,thick] (level4) -- (level5);
  
  \node[right=1cm of level1,text width=3cm] {\small ブルバキの一般理論};
  \node[right=1cm of level2,text width=3cm] {\small 可測性の階層};
  \node[right=1cm of level3,text width=3cm] {\small 情報の増加};
  \node[right=1cm of level4,text width=3cm] {\small 本稿の主題};
  \node[right=1cm of level5,text width=3cm] {\small 将来の展開};
\end{tikzpicture}
\end{center}

\subsection{Lean 4 コードへのリンク}

本稿の対象ファイル：

\begin{itemize}
\item \texttt{SigmaAlgebraTower\_Skeleton.lean}：$\sigma$-代数の構造塔
\item \texttt{StoppingTime\_MinLayer.lean}：停止時間と $\minLayer$ の関係
\end{itemize}

主要定義：

\begin{itemize}
\item SigmaAlgebraTower：71-108行目
\item フィルトレーションの構造塔化：263-286行目
\item 停止時間の定義：49-51行目
\item 停止過程：242-248行目
\end{itemize}

\subsection{今後の展望}

\begin{enumerate}
\item \textbf{測度の導入}：確率測度と積分の構造塔
\item \textbf{条件付き期待値}：$\minLayer$ を使った新しい特徴づけ
\item \textbf{マルチンゲール}：停止定理の構造塔的証明
\item \textbf{確率収束}：層の極限としての収束概念
\item \textbf{一般化}：連続時間や非可算添字への拡張
\end{enumerate}

\section*{参考文献}

\begin{enumerate}
\item Nicolas Bourbaki, \textit{Éléments de mathématique: Théorie des ensembles}
\item Patrick Billingsley, \textit{Probability and Measure}, Wiley (1995)
\item David Williams, \textit{Probability with Martingales}, Cambridge University Press (1991)
\item The mathlib Community, \textit{The Lean Mathematical Library}, \url{https://github.com/leanprover-community/mathlib4}
\item Rick Durrett, \textit{Probability: Theory and Examples}, Cambridge University Press (2019)
\end{enumerate}

\appendix

\section{Lean コードの詳細解説}

\subsection{SigmaAlgebraTower の実装}

完全な定義（71-108行目）：

\begin{lstlisting}[language=Lean,numbers=left,firstnumber=71]
def SigmaAlgebraTower : StructureTowerMin (Set Ω) (MeasurableSpace Ω) where
  layer 𝓕 := {A : Set Ω | @MeasurableSet Ω 𝓕 A}
  
  covering := by
    intro A
    use ⊤
    trivial
  
  monotone := by
    intro 𝓕₁ 𝓕₂ h𝓕 A hA
    exact h𝓕 A hA
  
  minLayer := fun A =>
    MeasurableSpace.generateFrom {A}
  
  minLayer_mem := by
    intro A
    exact MeasurableSpace.measurableSet_generateFrom
      (by simp : A ∈ ({A} : Set (Set Ω)))
  
  minLayer_minimal := by
    intro A 𝓕 hA
    apply MeasurableSpace.generateFrom_le
    intro B hB
    rcases hB with rfl
    exact hA
\end{lstlisting}

各フィールドの意味：

\begin{itemize}
\item \texttt{layer}：各 $\sigma$-代数に対応する可測集合の族
\item \texttt{covering}：全体 $\sigma$-代数 $\top$ を使って被覆性を示す
\item \texttt{monotone}：$\sigma$-代数の順序から層の包含を導く
\item \texttt{minLayer}：\texttt{generateFrom} を使って最小 $\sigma$-代数を構成
\item \texttt{minLayer\_mem}：生成 $\sigma$-代数の性質を使う
\item \texttt{minLayer\_minimal}：\texttt{generateFrom\_le} による最小性
\end{itemize}

\subsection{停止過程の実装}

停止過程の定義と主要な補題：

\begin{lstlisting}[language=Lean,numbers=left,firstnumber=242]
def stopped (X : ℕ → Ω → α) (τ : Ω → ℕ) (n : ℕ) : Ω → α :=
  fun ω => X (Nat.min n (τ ω)) ω

lemma stopped_eq_of_le (X : ℕ → Ω → α) (τ : Ω → ℕ)
    (n : ℕ) (ω : Ω) (hn : n ≤ τ ω) :
    stopped X τ n ω = X n ω := by
  simp [stopped, Nat.min_eq_left hn]

lemma stopped_const_after (X : ℕ → Ω → α) (τ : Ω → ℕ)
    (ω : Ω) (hτ : τ ω ≤ n) (hnm : n ≤ m) :
    stopped X τ m ω = stopped X τ n ω := by
  simp [stopped]
  have : Nat.min m (τ ω) = τ ω := Nat.min_eq_right (le_trans hτ hnm)
  have : Nat.min n (τ ω) = τ ω := Nat.min_eq_right hτ
  simp [*]
\end{lstlisting}

これらの補題により、停止過程の基本的な性質が証明される。

\end{document}
